\twocolumn[
\begin{@twocolumnfalse}
\begin{center}
{\Large\bfseries 大语言模型在轮椅转弯空间推理中的行为挖掘研究\par}
\vspace{0.25em}
{\normalsize Behavior Mining of Large Language Models in Wheelchair Turning Spatial Reasoning\par}
\vspace{0.45em}
{\small 单屹 \quad 张驰 \quad 徐敬研 \quad 傅玉 \quad 董梓苏 \quad 黄鸣辉\par}
\end{center}
\vspace{0.25em}
\noindent\textbf{摘要}\quad
本文研究大语言模型在轮椅转弯空间推理中的行为稳定性、噪声敏感性与失效模式。我们将“电动轮椅能否在楼梯转角处完成 $90^\circ$ 转弯”建模为一个受控文本推理任务，把场景表示为通道净宽、平台进深、轮椅宽度、轮椅长度和最小转弯半径等几何变量，并围绕描述层级、无关扰动和场景生成策略构造 Prompt 实验矩阵。与只报告单次准确率不同，本文把模型回复视为可挖掘的行为数据：每条记录同时保留 Prompt 条件、模型、重复编号、原始回复、解析判断和失败类别，从而分析模型何时依赖常识直觉，何时尝试形式化建模，以及哪些非因果线索会诱发判断翻转。正式实验选取 5 组代表性 Prompt 来源，每组 300 条，共 1500 条样本，并通过 DF/OpenAI-compatible 网关调用 12 个模型、每个样本重复采样 3 次，形成 $1500\times 12\times 3=54{,}000$ 条原始模型行为记录。本文给出任务定义、相关工作、Prompt 矩阵、多模型采集协议、行为特征接口和挖掘分析框架，作为完整实验数据分析与结果解释的基础。
\par\vspace{0.3em}
\noindent\textbf{关键词}\quad 大语言模型；空间推理；行为挖掘；轮椅转弯；扰动敏感性；失效模式
\vspace{0.85em}
\end{@twocolumnfalse}
]
